\documentclass[12pt]{article}

% Ћирилица: кодирање фонта и језик (српски, ћирилица)
\usepackage[T2A]{fontenc}
\usepackage[utf8]{inputenc}
\usepackage[serbianc]{babel}

% Величина стране и маргине
\usepackage[a4paper,top=2cm,bottom=2cm,left=2.5cm,right=2.5cm]{geometry}

% Корисни пакети
\usepackage{amsmath}
\usepackage{enumitem}
\usepackage{booktabs}
\usepackage{tcolorbox}
\tcbuselibrary{breakable}
\usepackage[colorlinks=true, allcolors=blue]{hyperref}

% Цртање електричних кола (европски симболи)
\usepackage[european]{circuitikz}
\usetikzlibrary{babel}

% Оквир за решене примере
\newtcolorbox[auto counter]{primer}{breakable,
    colback=white, colframe=black!55, colbacktitle=black!55, coltitle=white,
    fonttitle=\bfseries\small, title={Пример~\thetcbcounter},
    boxrule=0.5pt, arc=1.5pt, left=6pt, right=6pt, top=3pt, bottom=3pt,
    before skip=8pt, after skip=10pt}

% Оквир за кључне формуле
\newtcolorbox{kljucno}{colback=blue!4, colframe=blue!45!black,
    boxrule=0.5pt, arc=1.5pt, left=8pt, right=8pt, top=4pt, bottom=4pt}

\title{Везе отпорника и кондензатора}
\author{Основи електротехнике 1}
\date{}

\begin{document}
\maketitle
%==================================================================
\section{Проводност}
%==================================================================

Отпорност показује колико проводник \textbf{отежава} пролаз струје. Супротна
величина, која показује колико проводник \textbf{олакшава} пролаз струје,
зове се \textbf{проводност} и обележава се са $G$. Проводност је реципрочна
вредност отпорности:

\begin{kljucno}
\[
G = \frac{1}{R} \qquad [\text{S}]
\]
\end{kljucno}

Јединица за проводност је \textbf{сименс} (S), и важи
$1\,\text{S} = \dfrac{1}{1\,\Omega}$. Што је отпорност мања, проводност је
већа, и обрнуто.

Помоћу проводности Омов закон гласи:
\[
I = G \cdot U
\]

Код паралелне везе \textbf{проводности се сабирају}:
\[
G_e = G_1 + G_2 + \dots + G_n
\]
Зато формула за паралелну везу отпорника има реципрочне вредности:
$\dfrac{1}{R_e}$ није ништа друго него укупна проводност $G_e$.

Слично отпорности, и проводност проводника зависи од материјала. Уместо
специфичне отпорности $\rho$ тада користимо \textbf{специфичну проводност}
$\sigma = \dfrac{1}{\rho}$:
\[
G = \sigma \, \frac{S}{l}
\]

\begin{primer}
Израчунати проводност отпорника отпорности $R = 20\,\Omega$.

\textbf{Решење:}
\[
G = \frac{1}{R} = \frac{1}{20\,\Omega} = 0{,}05\,\text{S} = 50\,\text{mS}
\]
\end{primer}

\begin{primer}
Два отпорника, проводности $G_1 = 0{,}2\,\text{S}$ и $G_2 = 0{,}3\,\text{S}$,
везана су паралелно. Израчунати еквивалентну проводност и еквивалентну
отпорност.

\textbf{Решење:}

Код паралелне везе проводности се сабирају:
\[
G_e = G_1 + G_2 = 0{,}2\,\text{S} + 0{,}3\,\text{S} = 0{,}5\,\text{S}
\]
Еквивалентна отпорност је реципрочна вредност:
\[
R_e = \frac{1}{G_e} = \frac{1}{0{,}5\,\text{S}} = 2\,\Omega
\]
\end{primer}

%==================================================================
\section{Густина струје}
%==================================================================

Иста струја може да тече и кроз танак и кроз дебео проводник. У танком
проводнику струја је „збијена” на мању површину, па се он више загрева. Да
бисмо то описали, уводимо \textbf{густину струје} $J$. Она показује колика
струја пролази кроз јединицу површине попречног пресека проводника:

\begin{kljucno}
\[
J = \frac{I}{S} \qquad \left[\frac{\text{A}}{\text{mm}^2}\right]
\]
\end{kljucno}

У SI систему јединица за густину струје је $\text{A}/\text{m}^2$, али се у
пракси најчешће користи $\text{A}/\text{mm}^2$.

Што је густина струје већа, проводник се више загрева. Зато се за веће струје
користе проводници већег пресека: тако густина струје остаје у дозвољеним
границама и проводник се не прегрева.

\begin{primer}
Кроз проводник треба да тече струја $I = 12\,\text{A}$, а густина струје не
сме да пређе $J = 4\,\text{A}/\text{mm}^2$. Колики најмање мора бити пресек
проводника?

\textbf{Решење:}

Из формуле $J = \dfrac{I}{S}$ изражавамо пресек:
\[
S = \frac{I}{J} = \frac{12\,\text{A}}{4\,\text{A}/\text{mm}^2} = 3\,\text{mm}^2
\]
\end{primer}

%==================================================================
\section{Мешовита веза отпорника}
%==================================================================

Када у истом колу постоје и редна и паралелна веза, рачунамо \textbf{корак по
корак}: нађемо еквивалентну отпорност једне групе отпорника, заменимо целу
групу једним отпорником и понављамо поступак док не остане само један
отпорник.

\begin{primer}
Отпорници $R_2 = 6\,\Omega$ и $R_3 = 12\,\Omega$ везани су паралелно, а са њима
је редно везан отпорник $R_1 = 4\,\Omega$. Израчунати еквивалентну отпорност
између тачака $A$ и $B$.

\begin{center}
\begin{circuitikz}
    \draw (0,1) node[left]{$A$} to[short, o-] (0.5,1)
          to[R, l=$R_1$] (3,1) to[short, -*] (3.5,1);
    \draw (3.5,1) -- (3.5,2) to[R, l=$R_2$] (6.5,2) -- (6.5,1);
    \draw (3.5,1) -- (3.5,0) to[R, l=$R_3$] (6.5,0) -- (6.5,1);
    \draw (6.5,1) to[short, *-o] (7.5,1) node[right]{$B$};
\end{circuitikz}
\end{center}

\textbf{Решење:}

\textbf{Корак 1.} Отпорници $R_2$ и $R_3$ су паралелно, па их замењујемо
једним отпорником $R_{23}$:
\[
R_{23} = \frac{R_2 \cdot R_3}{R_2 + R_3}
       = \frac{6\,\Omega \cdot 12\,\Omega}{6\,\Omega + 12\,\Omega}
       = \frac{72}{18}\,\Omega = 4\,\Omega
\]

Сада коло изгледа овако:
\begin{center}
\begin{circuitikz}
    \draw (0,0) node[left]{$A$} to[short, o-] (0.5,0)
          to[R, l=$R_1$] (3,0)
          to[R, l=$R_{23}$] (5.5,0)
          to[short, -o] (6,0) node[right]{$B$};
\end{circuitikz}
\end{center}

\textbf{Корак 2.} Отпорници $R_1$ и $R_{23}$ су редно, па се сабирају:
\[
R_e = R_1 + R_{23} = 4\,\Omega + 4\,\Omega = 8\,\Omega
\]
\end{primer}


\subsection{Мешовита веза кондензатора}

Мешовита веза се решава исто као код отпорника, корак по корак.

\begin{primer}
Кондензатори $C_2 = 1\,\mu\text{F}$ и $C_3 = 2\,\mu\text{F}$ везани су
паралелно, а са њима је редно везан кондензатор $C_1 = 6\,\mu\text{F}$.
Израчунати еквивалентну капацитивност између тачака $A$ и $B$.

\begin{center}
\begin{circuitikz}
    \draw (0,1) node[left]{$A$} to[short, o-] (0.5,1)
          to[C, l=$C_1$] (3,1) to[short, -*] (3.5,1);
    \draw (3.5,1) -- (3.5,2) to[C, l=$C_2$] (6.5,2) -- (6.5,1);
    \draw (3.5,1) -- (3.5,0) to[C, l=$C_3$] (6.5,0) -- (6.5,1);
    \draw (6.5,1) to[short, *-o] (7.5,1) node[right]{$B$};
\end{circuitikz}
\end{center}

\textbf{Решење:}

\textbf{Корак 1.} Кондензатори $C_2$ и $C_3$ су паралелно, па се сабирају:
\[
C_{23} = C_2 + C_3 = 1\,\mu\text{F} + 2\,\mu\text{F} = 3\,\mu\text{F}
\]

\textbf{Корак 2.} Кондензатори $C_1$ и $C_{23}$ су редно:
\[
C_e = \frac{C_1 \cdot C_{23}}{C_1 + C_{23}}
    = \frac{6\,\mu\text{F} \cdot 3\,\mu\text{F}}{6\,\mu\text{F} + 3\,\mu\text{F}}
    = 2\,\mu\text{F}
\]
\end{primer}

\end{document}